\documentclass[../main.tex]{subfiles}
\begin{document}

Chương 2 trình bày tổng quan tài liệu và cơ sở lý thuyết làm nền tảng
cho việc nghiên cứu các giải pháp phát hiện bất thường dựa trên chuỗi
lời gọi hệ thống. Nội dung chương tập trung vào việc làm rõ phạm vi
nghiên cứu chung, phân tích các nghiên cứu liên quan trong lĩnh vực
phát hiện bất thường qua chuỗi lời gọi hệ thống (system calls), các
phương pháp trích xuất luật hoặc automata từ mô hình học sâu, và các
cơ chế giám sát thực thi thời gian thực bằng eBPF. Bên cạnh đó,
chương này cũng hệ thống hóa các kiến thức nền tảng về mạng
Transformer Decoder-only, lý thuyết Automata hữu hạn đơn định (DFA),
và kiến trúc eBPF trong nhân Linux, làm cơ sở khoa học để giải quyết
bài toán chưng cất tri thức từ mô hình học máy phức tạp thành luật
thực thi hiệu quả trong nhân hệ điều hành.

\section{Phạm vi nghiên cứu}
\label{sec:scope}

Trong bài toán phát hiện xâm nhập dựa trên lời gọi hệ thống (system
call) cho môi trường máy chủ, phạm vi nghiên cứu thường được giới hạn
và định nghĩa bởi các yếu tố cốt lõi bao gồm không gian giám sát, mô
hình đe dọa, và các bộ dữ liệu chuẩn được sử dụng để đánh giá.

Để đặt nền tảng cho việc thiết kế Guepard Shield, bài toán phát hiện
bất thường dựa trên syscall cần được hình thức hóa dưới góc độ của
mô hình hóa ngôn ngữ (language modeling) và lý thuyết automata.

Gọi $\mathcal{S} = \{s_1, s_2, \dots, s_{|V|}\}$ là tập hữu hạn các
tên lời gọi hệ thống (vocabulary). Một thực thi của tiến trình được
biểu diễn dưới dạng một chuỗi các sự kiện $X = (x_1, x_2, \dots, x_T)$,
trong đó mỗi $x_t \in \mathcal{S}$ là syscall tại thời điểm $t$.
Mục tiêu của pha huấn luyện mô hình Teacher (Transformer) là học một
hàm phân phối xác suất có điều kiện:
\begin{equation}
  P(x_t \mid x_{t-k}, \dots, x_{t-1})
\end{equation}
trong đó $k$ là chiều dài ngữ cảnh (context window). Một chuỗi quan
sát được coi là bất thường nếu xác suất xuất hiện của nó thấp hơn một
ngưỡng $\tau$ được định nghĩa trước. Tuy nhiên, việc tính toán trực
tiếp giá trị xác suất này trong nhân hệ điều hành là không khả thi do
các ràng buộc về tài nguyên đã nêu.

Thay vào đó, mục tiêu của luận văn là tìm một xấp xỉ rời rạc dưới dạng
Deterministic Finite Automaton (DFA), ký hiệu là $M = (Q, \Sigma,
\delta, q_0, F)$.
Trong đó:
\begin{itemize}
  \item $Q$ là tập hữu hạn các trạng thái, được trích xuất từ việc
    phân cụm không gian vector ẩn của Transformer.
  \item $\Sigma \subseteq \mathcal{S}$ là bảng chữ cái (tập syscall).
  \item $\delta: Q \times \Sigma \to Q$ là hàm chuyển trạng thái.
  \item $q_0 \in Q$ là trạng thái bắt đầu.
  \item $F \subset Q$ là tập các trạng thái chấp nhận (hành vi bình thường).
\end{itemize}
Thách thức của việc hình thức hóa này nằm ở chỗ Transformer có khả năng
ghi nhớ ngữ cảnh dài thông qua cơ chế chú ý, trong khi DFA truyền thống
chỉ phụ thuộc vào trạng thái hiện tại. Luận văn cần giải quyết mâu thuẫn
này bằng cách chọn lựa một biểu diễn trạng thái ẩn đủ giàu thông tin để
duy trì tính "Markov" cần thiết cho DFA mà không làm bùng nổ số lượng
trạng thái.

Một vấn đề nghiên cứu then chốt khác là khả năng chống lại các cuộc
tấn công bắt chước (Mimicry Attacks). Trong bối cảnh syscall, kẻ tấn
công có thể chèn thêm các lời gọi hệ thống "vô hại" (no-op hoặc các
syscall bình thường) vào giữa các bước của mã độc để làm thay đổi
biểu diễn vector của chuỗi, từ đó đánh lừa các hệ thống phát hiện
bất thường. Việc trích xuất DFA cần đảm bảo rằng cấu trúc whitelist
của automata đủ chặt chẽ để các hành vi chèn đệm này vẫn đưa con trỏ
trạng thái vào vùng tấn công, hoặc ít nhất là đưa vào các trạng thái
nghi vấn (suspect states) để kích hoạt cơ chế kiểm tra sâu hơn ở không gian
người dùng. Khả năng chịu đựng của DFA đối với các biến đổi trình tự
này là một trong những trọng tâm phân tích của luận văn.

Về không gian giám sát, các nghiên cứu trong lĩnh vực này tập trung
vào bài toán phát hiện bất thường ở mức lời gọi hệ thống
(syscall-level anomaly detection) trên các máy chủ Linux chạy các
khối lượng công việc container hóa (containerized workloads). Giao
diện system call đại diện cho ranh giới phân tách thấp nhất giữa
không gian người dùng (user space) và không gian nhân (kernel space).
Mọi tương tác của ứng dụng với tài nguyên phần cứng, mạng, tệp tin,
hoặc tiến trình khác đều bắt buộc phải đi qua ranh giới này. Do đó,
việc giám sát ở mức lời gọi hệ thống cung cấp một cái nhìn toàn diện
và khó bị vượt qua bởi các hành vi xâm nhập thông thường.

Mô hình đe dọa mục tiêu thường tập trung vào giai đoạn hậu xâm nhập
(post-exploitation), giả định rằng kẻ tấn công đã thực thi thành công
mã độc hoặc chiếm được quyền truy cập vào container thông qua các lỗ
hổng phần mềm. Các hành vi xâm nhập mục tiêu bao gồm thiết lập kết
nối ngược (reverse shell), leo thang đặc quyền (privilege
escalation), đánh cắp dữ liệu (data exfiltration), và duy trì sự hiện
diện lâu dài trong hệ thống (persistence). Kẻ tấn công được coi là
hoạt động dưới dạng hộp đen (black-box), nghĩa là họ không có tri
thức về các cơ chế chính sách bảo mật đang được thực thi tại runtime
để giám sát hành vi của hệ thống.

Phạm vi nghiên cứu này không giải quyết các vấn đề liên quan đến tấn
công né tránh đối kháng (adversarial evasion) nhắm trực tiếp vào việc
thay đổi phân phối chuỗi syscall để đánh lừa mô hình học máy. Các
rootkit mức kernel can thiệp trực tiếp vào cấu trúc dữ liệu của hệ
điều hành hoặc các hệ thống phát hiện xâm nhập mức mạng (Network IDS)
cũng nằm ngoài phạm vi khảo sát của luận văn.

\section{Các nghiên cứu liên quan}
\label{sec:related_work}

\subsection{Phát hiện bất thường dựa trên chuỗi lời gọi hệ thống}
Các nghiên cứu về hệ thống phát hiện xâm nhập trên máy chủ (HIDS) dựa
trên syscall đã trải qua nhiều giai đoạn phát triển, từ các thuật
toán thống kê đơn giản đến các mô hình học sâu phức tạp.

Phương pháp kinh điển đầu tiên là STIDE (Sequence Time-Delay
Embedding) do Forrest và các cộng sự đề xuất \cite{forrest1996sense}.
STIDE hoạt động dựa trên cơ chế n-gram, xây dựng một cơ sở dữ liệu
chứa tất cả các chuỗi lời gọi hệ thống bình thường có độ dài cố định
$n$ thu được trong giai đoạn huấn luyện. Khi vận hành thực tế, bất kỳ
chuỗi syscall con nào có độ dài $n$ xuất hiện mà không nằm trong cơ
sở dữ liệu bình thường sẽ bị coi là bất thường. Mặc dù STIDE có tốc
độ thực thi rất nhanh và thuật toán đơn giản, mô hình này gặp hạn chế
lớn do bỏ lỡ các phụ thuộc ngữ cảnh dài hạn và dễ sinh cảnh báo giả
khi ứng dụng có nhiều luồng thực thi hợp lệ xen kẽ. Các cải tiến sau
đó sử dụng Mô hình Markov ẩn (Hidden Markov Model - HMM) để mô hình
hóa sự chuyển dịch trạng thái của tiến trình dưới dạng xác suất
\cite{warrender1999detecting}, tuy
nhiên chi phí tính toán tăng rất nhanh khi kích thước từ vựng
(vocabulary) của các syscall mở rộng.

Sự trỗi dậy của học sâu đã thay đổi căn bản cách tiếp cận này.
DeepLog \cite{du2017deeplog} giới thiệu việc áp dụng mạng nơ-ron hồi
quy (RNN/LSTM) \cite{hochreiter1997lstm} cho bài toán phát hiện bất thường thông qua việc dự
đoán phần tử tiếp theo (next-token prediction) trên chuỗi log hoặc
chuỗi syscall. Khi xác suất dự đoán của syscall tiếp theo thực tế
thấp hơn một ngưỡng định sẵn, hệ thống sẽ gắn cờ cảnh báo. Để giải
quyết hiện tượng rò rỉ dữ liệu khi trích xuất đặc trưng bằng Word2Vec
\cite{mikolov2013word2vec} truyền thống trên chuỗi syscall, Mvula và các cộng sự
\cite{mvula2023evaluating} khuyến nghị sử dụng các phương pháp mã hóa
nén như Byte-Pair Encoding (BPE) \cite{sennrich2016bpe}. Dù đạt độ chính xác cao, các mô
hình dựa trên LSTM gặp khó khăn lớn trong việc xử lý song song khi
huấn luyện và có độ trễ suy luận lớn, không đáp ứng được yêu cầu phản
hồi thời gian thực trong nhân hệ điều hành.

Gần đây, cơ chế Self-Attention trong kiến trúc Transformer đã được
ứng dụng để mô hình hóa chuỗi syscall. Nghiên cứu của Fournier và các
cộng sự \cite{fournier2023language} đã chỉ ra rằng các mô hình ngôn
ngữ lớn (như Transformer và Longformer) đạt độ chính xác F1 và AUROC
vượt trội (>95\%) trên các bài toán phát hiện bất thường nhờ khả năng
nắm bắt tốt các quan hệ ngữ cảnh dài hạn giữa các syscall đứng xa
nhau. TransCall \cite{alsiam2026transcall} cũng khai thác Transformer
để phát hiện mã độc zero-day dưới dạng xử lý chuỗi syscall tương tự
như ngôn ngữ tự nhiên. Tuy nhiên, rào cản lớn nhất của các mô hình
Transformer là chi phí tính toán và bộ nhớ cực kỳ lớn. Việc chạy trực
tiếp một mô hình Transformer trong nhân hệ điều hành là hoàn toàn bất
khả thi do nhân Linux không hỗ trợ các phép tính dấu phẩy động và yêu
cầu độ trễ xử lý của mỗi syscall phải ở mức micro giây. Khoảng cách
này tạo ra một "khoảng cách thực thi" (enforcement gap) lớn giữa mô
hình học máy và môi trường runtime thực tế.

\subsection{Trích xuất luật khả diễn và Automata từ mạng nơ-ron}
Để vượt qua tính chất "hộp đen" của các mô hình học sâu và đảm bảo
tính giải trình trong các hệ thống an ninh mạng (forensic
accountability), hướng nghiên cứu về mô hình thế thân giải thích được
(interpretable surrogate models) và chưng cất tri thức (Knowledge
Distillation) đã nhận được nhiều sự quan tâm. Thay vì cố gắng giải
thích các mô hình phức tạp thông qua các công cụ hậu kiểm (post-hoc)
như LIME \cite{ribeiro2016lime} hay SHAP \cite{lundberg2017shap} vốn tốn thời gian suy luận phụ và không đảm bảo
tính nhất quán, hướng tiếp cận này tìm cách chuyển đổi tri thức từ
mạng nơ-ron sang các biểu diễn tường minh, dễ hiểu và có thể thực thi hiệu quả.

\subsubsection{Cơ sở lý thuyết về chưng cất tri thức (Knowledge Distillation)}
Chưng cất tri thức \cite{hinton2015distilling} là một kỹ thuật trong đó một mô hình nhỏ hơn, gọn
hơn (Student) được huấn luyện để mô phỏng hành vi của một mô hình
lớn, phức tạp hơn (Teacher). Trong bài toán phát hiện xâm nhập, mô
hình Teacher (như Transformer) đóng vai trò là một bộ phân loại chính
xác nhưng nặng nề, trong khi Student (như DFA) cần đạt được độ trung
thực (fidelity) cao nhất có thể so với Teacher tại các vùng quyết
định nhạy cảm. Quá trình này không chỉ là việc bắt chước nhãn đầu ra
(hard labels) mà còn là việc học theo phân phối xác suất (soft
targets) của Teacher, giúp Student nắm bắt được ranh giới quyết định
mượt mà hơn.

\subsubsection{Kỹ thuật trích xuất Automata từ RNN và Transformer}
Trong số các biểu diễn rời rạc, Automata hữu hạn đơn định (DFA) nổi
lên như một ứng viên sáng giá nhờ tính đơn định và chi phí thực thi
$O(1)$. Công trình tiêu biểu của Weiss và các cộng sự
\cite{weiss2018extracting} đề xuất kỹ thuật trích xuất DFA từ các mạng
hồi quy (RNNs) thông qua quy trình gom cụm không gian trạng thái ẩn.
Cụ thể, các vector trạng thái ẩn $h_t$ (vốn mang thông tin lịch sử
của chuỗi đầu vào) được ánh xạ vào một tập hữu hạn các trạng thái rời
rạc $Q$ bằng các thuật toán phân cụm như K-Means hoặc Mean-Shift.

Đối với kiến trúc Transformer, thách thức nằm ở cơ chế self-attention
vốn không duy trì một trạng thái ẩn tuần tự duy nhất như RNN. Tuy
nhiên, theo nguyên lý tương tự kỹ thuật của Weiss và các cộng sự
\cite{weiss2018extracting}, vector nhúng lớp cuối cùng (final-layer
embeddings) của Transformer vẫn mã hóa đủ thông tin ngữ cảnh nhân
quả để có thể rời rạc hóa bằng phân cụm. Nghiên cứu của Herbinger và
các cộng sự \cite{herbinger2023leveraging} về mô hình thế thân giải
thích được (interpretable surrogate models) củng cố thêm rằng biểu
diễn rời rạc học được từ đặc trưng nội tại của mô hình phức tạp có
thể xấp xỉ trung thực hành vi của Teacher. Việc trích xuất DFA từ Transformer yêu
cầu một bước xử lý bổ sung để giải quyết tính không đơn định khi ánh
xạ từ không gian liên tục sang các bước chuyển trạng thái rời rạc,
thường thông qua các chiến lược cắt tỉa thống kê (statistical
pruning) để loại bỏ các chuyển dịch nhiễu hoặc có xác suất thấp.

\subsubsection{Ứng dụng trong an ninh hệ thống và pháp y kỹ thuật số}
Việc chuyển đổi tri thức sang DFA mang lại lợi ích kép cho an ninh hệ
thống. Về mặt vận hành, DFA cho phép triển khai cơ chế whitelist động
có khả năng phản hồi ở tốc độ của phần cứng hoặc kernel. Về mặt phân
tích, cấu trúc đồ thị của DFA cho phép các chuyên gia bảo mật truy
vết ngược lại các lộ trình tấn công dưới dạng các chuỗi trạng thái
tường minh. Nếu một chuỗi syscall bị gắn attack, DFA có thể chỉ ra
chính xác syscall nào và tại ngữ cảnh nào đã gây ra sự vi phạm, điều
mà các mô hình Transformer nguyên bản không thể cung cấp một cách trực quan.

Tuy nhiên, khoảng trống nghiên cứu chính nằm ở việc tối ưu hóa quy
trình chưng cất này cho các ràng buộc của môi trường thực thi kernel
như eBPF. Hầu hết các công trình hiện nay mới chỉ tập trung vào việc
đạt được độ trung thực cao về mặt toán học mà chưa xem xét đến giới
hạn bộ nhớ của BPF Maps hoặc khả năng xử lý các chuỗi syscall có
chiều dài biến thiên trong môi trường máy chủ thực tế. Guepard Shield
hướng tới việc lấp đầy khoảng trống này bằng một pipeline tự động hóa
từ trích xuất trạng thái ẩn của Transformer Decoder-only đến nạp bảng
chuyển DFA vào eBPF.

\subsection{Thực thi chính sách bảo mật thời gian thực qua eBPF}
Công nghệ eBPF (Extended Berkeley Packet Filter) đại diện cho một
bước tiến quan trọng trong kiến trúc nhân Linux, cho phép chạy các
chương trình sandboxed hiệu năng cao trực tiếp tại các điểm móc
(hooks) như tracepoints hoặc LSM (Linux Security Modules) với chi phí
chuyển đổi ngữ cảnh (context switch) bằng không.

Các công cụ bảo mật cloud-native phổ biến hiện nay như Falco và
Tetragon \cite{her2025indepth} tận dụng eBPF để thu thập telemetry và
thực thi chính sách ở mức nhân. Nghiên cứu so sánh hiệu năng của Her
và các cộng sự \cite{her2025indepth} cho thấy Tetragon có thể phát
hiện và ngăn chặn các hành vi thoát khỏi container (container
escapes) chỉ trong 1.178 giây với mức chiếm dụng CPU tối ưu (~73\%)
so với các giải pháp dựa trên không gian người dùng truyền thống. Tuy
nhiên, điểm hạn chế cốt lõi của các công cụ này là toàn bộ tập luật
(ruleset) phát hiện đều phải được viết thủ công bằng YAML dựa trên
tri thức của chuyên gia bảo mật. Cách tiếp cận này mang tính tĩnh,
khó bao phủ các chuỗi hành vi phức tạp và dễ bị vượt qua bởi các biến
thể tấn công zero-day.

Để khắc phục điều này, một số nghiên cứu gần đây tìm cách đưa mô hình
học máy vào thực thi trong kernel qua eBPF. Zhang và các cộng sự
\cite{zhang2024realtime} đã thiết kế và triển khai một mạng nơ-ron
nhân tạo trực tiếp trong kernel bằng eBPF bằng cách lượng tử hóa mô
hình sang số học số nguyên (integer-only arithmetic) để vượt qua giới
hạn stack 512 bytes của eBPF verifier. Mô hình đạt thời gian suy luận
cực nhanh (3-5 micro giây) nhưng chương trình eBPF này cực kỳ phức
tạp, khó bảo trì, và mô hình nơ-ron chạy trong kernel vẫn là một hộp
đen không thể giải thích trực tiếp. Bachl và các cộng sự
\cite{bachl2021flow} cũng chứng minh việc thực thi cây quyết định
(decision tree) trong eBPF tăng 20\% hiệu suất so với không gian người dùng.

Bảng~\ref{tab:topt} tóm tắt các ưu điểm và hạn chế chính của các hướng
tiếp cận liên quan.

\begin{table}[htbp]
  \caption{So sánh các phương pháp phát hiện xâm nhập dựa trên syscall}
  \label{tab:topt}
  \centering
  \resizebox{\textwidth}{!}{
    \begin{tabular}{|l|p{3.5cm}|p{3.5cm}|p{2cm}|p{2cm}|}
      \hline
      \textbf{Phương pháp} & \textbf{Ưu điểm} & \textbf{Hạn chế} &
      \textbf{Khả năng học dữ liệu} & \textbf{Thực thi trong Kernel} \\ \hline
      \textbf{STIDE / N-gram} \cite{forrest1996sense} & Tốc độ nhanh,
      thuật toán đơn giản, dễ cài đặt. & Bỏ lỡ các ngữ cảnh dài hạn,
      độ chính xác thấp trên hệ thống phức tạp. & Có (Thống kê tần
      suất) & Có (Thông qua cấu trúc seccomp-BPF thô) \\ \hline
      \textbf{LSTM / DeepLog} \cite{du2017deeplog} & Nắm bắt ngữ cảnh
      động tốt hơn n-gram, độ chính xác khá tốt. & Khó xử lý song
      song khi huấn luyện, độ trễ suy luận lớn. & Có (Học sâu có giám
      sát/không giám sát) & Không phù hợp để chạy trực tiếp trong
      kernel \\ \hline
      \textbf{Transformer HIDS} \cite{fournier2023language,
      alsiam2026transcall} & Mô hình hóa tốt phụ thuộc ngữ cảnh dài,
      thường đạt kết quả phát hiện cao. & Chi phí suy luận và bộ nhớ
      lớn, khó đặt trực tiếp trên đường đi syscall. & Có (Cơ chế tự
      chú ý causal) & Không phù hợp để chạy trực tiếp trong kernel \\ \hline
      \textbf{Rules (Falco / Tetragon)} \cite{her2025indepth} & Chạy
      trực tiếp trong kernel với eBPF, độ trễ thấp. & Luật viết thủ
      công bởi chuyên gia, khó bao phủ biến thể trình tự phức tạp. &
      Không có (Tĩnh) & Có (Được thiết kế chạy trực tiếp bằng eBPF) \\ \hline
      \textbf{Guepard Shield (Đề xuất)} & Chuyển tri thức từ
      Transformer sang biểu diễn trạng thái hữu hạn để kiểm tra nhanh
      bằng eBPF. & Phụ thuộc vào chất lượng mô hình Teacher và cách
      rời rạc hóa trạng thái. & Có (Chưng cất từ Transformer Teacher)
      & Có (Thực thi bảng chuyển trạng thái
      DFA bằng eBPF Maps) \\ \hline
    \end{tabular}
  }
\end{table}

Từ các phân tích trên, khoảng trống lớn nhất là việc tự động hóa cầu
nối giữa mô hình học sâu tuần tự (Transformer) và bộ thực thi chính
sách eBPF. Chưa có giải pháp nào tự động chưng cất tri thức từ mô
hình Transformer thành một bảng chuyển trạng thái DFA tối ưu, sau đó
nạp bảng chuyển này vào các cấu trúc BPF Maps để thực thi giám sát
với độ phức tạp $O(1)$ và khả năng giải thích hoàn toàn. Đây là một
khoảng trống kỹ thuật lớn trong việc kết nối khả năng học ngữ cảnh
của Transformer với hiệu năng thực thi của eBPF.

\section{Mạng Transformer Decoder-only}
\label{sec:transformer}

Kiến trúc Transformer Decoder-only \cite{vaswani2017attention} là mô hình mạng nơ-ron tự hồi quy
(autoregressive) được ứng dụng rộng rãi để xử lý và dự đoán các chuỗi
dữ liệu tuần tự. Kiến trúc này loại bỏ phần Encoder và chỉ sử dụng
các lớp causal self-attention để mô hình hóa luồng sự kiện theo trình
tự thời gian.

Giả sử đầu vào của mô hình là một chuỗi các token biểu diễn các lời
gọi hệ thống $S = (s_1, s_2, \ldots, s_T)$, trong đó mỗi $s_t$ thuộc
một từ vựng hữu hạn $\Sigma$ đại diện cho tập hợp các syscall được
định danh. Chuỗi token này trước tiên được đi qua một lớp nhúng
(embedding layer) để chuyển thành chuỗi các vector biểu diễn $H^{(0)}
= (h_1^{(0)}, h_2^{(0)}, \ldots, h_T^{(0)})$ với $h_t^{(0)} \in \mathbb{R}^d$.

Tại mỗi lớp self-attention thứ $l$, mô hình tính toán các ma trận
Query ($Q$), Key ($K$) và Value ($V$) từ biểu diễn ẩn của lớp trước
đó $H^{(l-1)}$ bằng các trọng số tương ứng $W_Q^{(l)}, W_K^{(l)},
W_V^{(l)} \in \mathbb{R}^{d \times d}$:
\begin{equation}
  Q^{(l)} = H^{(l-1)} W_Q^{(l)}, \quad K^{(l)} = H^{(l-1)} W_K^{(l)},
  \quad V^{(l)} = H^{(l-1)} W_V^{(l)}
\end{equation}

Để đảm bảo tính nhân quả (causality) trong việc dự đoán phần tử tiếp
theo - nghĩa là mô hình tại thời điểm $t$ chỉ được phép chú ý đến các
phần tử đứng trước $s_1, \ldots, s_t$ mà không được nhìn thấy tương
lai $s_{t+1}, \ldots, s_T$ - một ma trận mặt nạ (causal mask) $M \in
\mathbb{R}^{T \times T}$ được cộng vào trước khi tính toán hàm
softmax. Các phần tử của ma trận $M$ được định nghĩa như sau:
\begin{equation}
  M_{i,j} =
  \begin{cases}
    0 & \text{nếu } i \ge j \\
    -\infty & \text{nếu } i < j
  \end{cases}
\end{equation}

Công thức tính toán của lớp causal self-attention được viết dưới dạng:
\begin{equation}
  \text{Attention}(Q^{(l)}, K^{(l)}, V^{(l)}) =
  \text{softmax}\left(\frac{Q^{(l)} (K^{(l)})^T}{\sqrt{d}} + M\right) V^{(l)}
\end{equation}

Mô hình được huấn luyện theo phương pháp học không giám sát bằng cách
tối thiểu hóa hàm mất mát entropy chéo cho bài toán dự đoán token
tiếp theo (next-token prediction) trên tập dữ liệu chuỗi syscall lành tính:
\begin{equation}
  \label{eq:loss}
  \mathcal{L} = -\frac{1}{T-1} \sum_{t=1}^{T-1} \log P(s_{t+1} \mid
  s_1, s_2, \ldots, s_t)
\end{equation}

Trong đó, xác suất điều kiện được tính thông qua lớp tuyến tính cuối
cùng (projection layer) và hàm softmax trên vector trạng thái ẩn lớp
cuối $h_t^{(L)} \in \mathbb{R}^d$:
\begin{equation}
  P(s_{t+1} = \sigma \mid s_1, \ldots, s_t) = \frac{\exp(w_\sigma^T
  h_t^{(L)} + b_\sigma)}{\sum_{\sigma' \in \Sigma} \exp(w_{\sigma'}^T
  h_t^{(L)} + b_{\sigma'})}
\end{equation}

Vector trạng thái ẩn $h_t^{(L)}$ tại vị trí $t$ tích lũy toàn bộ
thông tin lịch sử của luồng syscall từ bước $1$ đến $t$. Đặc tính này
làm cho vector trạng thái ẩn lớp cuối cùng trở thành đối tượng nghiên
cứu tiềm năng trong các kỹ thuật phân cụm trạng thái liên tục để
trích xuất automata hữu hạn rời rạc phục vụ mục đích giải thích hoặc
thực thi chính sách.

\section{Automata hữu hạn đơn định}
\label{sec:dfa}

Một Automata hữu hạn đơn định (Deterministic Finite Automaton - DFA)
là một mô hình toán học biểu diễn một hệ thống có số lượng trạng thái
hữu hạn, dịch chuyển giữa các trạng thái dựa trên các ký hiệu đầu vào
đơn định. Về mặt hình thức, một DFA được định nghĩa là một bộ 5 thành phần:
\begin{equation}
  A = (Q, \Sigma, \delta, q_0, F)
\end{equation}
Trong đó:
\begin{itemize}
  \item $Q$ là tập hợp hữu hạn và không rỗng các trạng thái ($Q =
    \{q_0, q_1, \ldots, q_m\}$).
  \item $\Sigma$ là bảng chữ cái đầu vào hữu hạn (trong bài toán
      syscall là tập hợp các chỉ số hoặc tên của lời gọi hệ thống nằm
    trong Vocabulary).
  \item $\delta: Q \times \Sigma \to Q$ là hàm chuyển trạng thái đơn
    định. Với mỗi trạng thái hiện tại $q \in Q$ và ký hiệu đầu vào $s
    \in \Sigma$, tồn tại duy nhất một trạng thái tiếp theo $q' =
    \delta(q, s) \in Q$.
  \item $q_0 \in Q$ là trạng thái khởi đầu của hệ thống.
  \item $F \subseteq Q$ là tập hợp các trạng thái chấp nhận (accepting states).
\end{itemize}

Trong bài toán phát hiện bất thường sử dụng DFA, tập các trạng thái
chấp nhận $F$ thường được thay thế bằng tập các trạng thái tấn công
(attack hoặc abnormal states) $A \subseteq Q$. Khi chuỗi syscall dẫn
DFA đến một trạng thái thuộc $A$, hệ thống sẽ phát tín hiệu cảnh báo xâm nhập.

Đặc tính quan trọng nhất khiến DFA trở thành mục tiêu triển khai lý
tưởng trong các hệ thống giám sát hiệu năng cao là tính hiệu quả cực
cao về mặt tính toán. Việc chuyển trạng thái chỉ yêu cầu một phép tra
cứu bảng chuyển dịch (transition lookup) với độ phức tạp thời gian là
$O(1)$ cho mỗi ký hiệu đầu vào:
\begin{equation}
  \text{Thời gian xử lý} = O(1), \quad \text{Bộ nhớ lưu trữ} \propto
  |Q| \times |\Sigma|
\end{equation}
DFA không yêu cầu bất kỳ phép tính số thực dấu phẩy động nào và có
thể được biểu diễn hoàn toàn bằng các cấu trúc dữ liệu mảng hoặc bảng
băm số nguyên tĩnh, hoàn toàn tương thích với các ràng buộc khắt khe
của môi trường nhân hệ điều hành.

\section{Công nghệ eBPF}
\label{sec:ebpf}

eBPF là một công nghệ mang tính cách mạng trong nhân Linux, cho phép
chạy các chương trình tùy biến trực tiếp trong không gian nhân
(kernel space) khi các sự kiện hệ thống xảy ra mà không cần sửa đổi
mã nguồn kernel hoặc tải các kernel module.

\subsection{Sự tiến hóa từ Classic BPF đến eBPF}
Nguyên thủy, BPF (Classic BPF - cBPF) được giới thiệu vào năm 1992
\cite{mccanne1993bpf} như một cơ chế lọc gói tin hiệu quả trong nhân. Tuy nhiên, kể từ
phiên bản nhân Linux 3.15, eBPF đã mở rộng phạm vi ứng dụng từ mạng
sang giám sát toàn diện hệ thống. Sự khác biệt cốt lõi nằm ở việc
eBPF sở hữu tập lệnh rộng hơn (64-bit thay vì 32-bit), số lượng thanh
ghi tăng từ 2 lên 10, và khả năng gọi các hàm helper trong nhân. Đặc
biệt, eBPF giới thiệu cơ chế JIT (Just-In-Time) compiler, cho phép
biên dịch mã bytecode eBPF trực tiếp sang mã máy bản địa của CPU,
giúp chương trình đạt hiệu năng thực thi gần tương đương với mã nguồn
được biên dịch trực tiếp vào nhân.

\subsection{Cơ chế hoạt động và Điểm móc (Hooks)}
Các chương trình eBPF hoạt động theo cơ chế hướng sự kiện
(event-driven). Chúng được gắn (hooked) vào các vị trí thực thi cụ
thể trong kernel, chẳng hạn như tracepoints (ví dụ:
  \texttt{sys\_enter} kích hoạt khi tiến trình chuẩn bị thực hiện một
syscall), kprobes/kretprobes (móc vào các hàm nhân bất kỳ), hoặc các
LSM hooks (Linux Security Modules) như AppArmor hay SELinux. Khi sự
kiện xảy ra, chương trình eBPF được thực thi để phân tích ngữ cảnh sự
kiện đó. Khả năng truy cập vào cấu trúc dữ liệu của nhân thông qua
các điểm móc này cho phép eBPF quan sát được các thông tin nhạy cảm
như PID, UID, tham số syscall, và ngữ cảnh tiến trình cha-con.

\subsection{Trình xác thực eBPF (Verifier) và Tính an toàn}
Để đảm bảo an toàn tuyệt đối cho hệ điều hành, nhân Linux tích hợp
một trình xác thực tĩnh (verifier) thực hiện kiểm tra mã nguồn eBPF
trước khi nạp vào nhân. Trình xác thực áp đặt các ràng buộc an toàn
vô cùng nghiêm ngặt:
\begin{itemize}
  \item \textbf{Giới hạn ngăn xếp (Stack limit):} Kích thước stack
    tối đa của một chương trình eBPF chỉ là 512 bytes. Mọi biến cục
    bộ và cấu trúc dữ liệu tạm thời bắt buộc phải nằm trong giới hạn
    này. Điều này ngăn chặn việc tiêu tốn quá mức bộ nhớ stack của
    nhân (kernel stack exhaustion).
  \item \textbf{Đảm bảo kết thúc (Termination guarantees):} Trình xác
    thực kiểm tra đồ thị dòng chảy điều khiển (Control Flow Graph -
    CFG) để đảm bảo không có vòng lặp vô hạn và số lượng chỉ thị lệnh
    tối đa không vượt quá giới hạn cho phép (thường là 1 triệu chỉ
    thị lệnh). Mọi nhánh rẽ trong chương trình phải được chứng minh
    là sẽ kết thúc.
  \item \textbf{Truy cập bộ nhớ an toàn:} Cấm truy cập trực tiếp vào
    các vùng nhớ nhân ngẫu nhiên; mọi truy cập bộ nhớ phải qua các
    hàm helper được kiểm duyệt (như \texttt{bpf\_probe\_read\_kernel}).
\end{itemize}

\subsection{Cấu trúc BPF Maps và Lưu trữ trạng thái}
Vì kích thước stack bị giới hạn ở 512 bytes, các chương trình eBPF
không thể lưu trữ các cấu trúc dữ liệu lớn trực tiếp trên stack. Thay
vào đó, eBPF cung cấp cấu trúc BPF Maps. Đây là các kho lưu trữ dữ
liệu dạng khóa - giá trị (key-value stores) hiệu năng cao nằm trong
bộ nhớ nhân, hỗ trợ tra cứu với độ phức tạp $O(1)$. Các loại BPF Maps
phổ biến bao gồm: \texttt{BPF\_MAP\_TYPE\_HASH},
\texttt{BPF\_MAP\_TYPE\_ARRAY}, hoặc \texttt{BPF\_MAP\_TYPE\_LRU\_HASH}.

Bảng chuyển trạng thái của DFA $\delta(q, s) = q'$ có thể được biểu
diễn trong kernel bằng một Hash Map của BPF. Trong cấu trúc này, khóa
(key) là một cặp số nguyên bao gồm định danh trạng thái hiện tại và
mã số lời gọi hệ thống (syscall ID), còn giá trị (value) tương ứng là
trạng thái kế tiếp. Khi một syscall được kích hoạt, chương trình eBPF
thực hiện phép tra cứu $O(1)$ trên Map này để cập nhật trạng thái DFA
của luồng tiến trình hiện tại (thường được lưu trong một
\texttt{per-thread state map} riêng biệt).

\subsection{Tính di động và eBPF Ecosystem (BTF \& CO-RE)}
Thách thức lớn nhất của eBPF truyền thống là sự phụ thuộc vào cấu
trúc dữ liệu cụ thể của từng phiên bản nhân Linux (kernel version
dependency). Để giải quyết vấn đề này, công nghệ BTF (BPF Type
Format) và cơ chế CO-RE (Compile Once - Run Everywhere) đã được giới
thiệu. BTF cung cấp siêu dữ liệu (metadata) về các kiểu dữ liệu trong
nhân, trong khi CO-RE cho phép trình nạp (loader) ở không gian người
dùng thực hiện việc điều chỉnh địa chỉ (relocation) mã bytecode eBPF
tại thời điểm chạy. Nhờ đó, một chương trình eBPF có thể chạy trên
nhiều phiên bản nhân khác nhau mà không cần biên dịch lại, một yếu tố
then chốt cho việc triển khai Guepard Shield trên các môi trường đám
mây đa dạng.

\subsection{Cập nhật trực tiếp (Atomic Update \& Hot-reload)}
Một đặc tính ưu việt của BPF Maps là hỗ trợ các thao tác cập nhật
nguyên tử (atomic map updates). Điều này cho phép một tiến trình chạy
ở không gian người dùng có thể ghi đè bảng chuyển trạng thái DFA mới
vào BPF Maps mà không cần dừng hoặc dỡ bỏ chương trình eBPF đang giám
sát trong kernel. Khả năng này đảm bảo hệ thống có thể cập nhật động
các chính sách bảo mật hoặc điều chỉnh trạng thái DFA khi phát hiện
sự dịch chuyển phân phối (concept drift) hoặc khi cần tối ưu hóa
automata dựa trên dữ liệu mới mà không gây gián đoạn quá trình giám
sát thời gian thực.

Tóm lại, công nghệ eBPF không chỉ cung cấp một hạ tầng thu thập dữ
liệu hiệu năng cao mà còn đóng vai trò là một bộ thực thi chính sách
(enforcement engine) mạnh mẽ. Các ràng buộc của verifier và đặc tính
của BPF Maps là những yếu tố quyết định đến cách chúng ta thiết kế và
rời rạc hóa mô hình DFA từ Transformer. Việc tận dụng BTF và CO-RE
giúp Guepard Shield đạt được tính linh hoạt và khả năng triển khai
rộng rãi trong các hệ thống Linux hiện đại.

\end{document}
